% ============================================================================
%  CAPÍTULO IV: MEDIDAS DE MITIGACIÓN Y ENDURECIMIENTO TÉCNICO
% ============================================================================

\section{Estrategia Integral de Remediación y Endurecimiento}

Para neutralizar de forma definitiva las vulnerabilidades identificadas durante el escaneo de Nmap NSE, se diseñó e implementó un plan técnico de mitigación fundamentado en el estándar de defensa en profundidad (\textit{Defense in Depth}) y las guías de endurecimiento del CIS (\textit{Center for Internet Security}) \citep{cis_debian_2023}. A continuación se detallan las medidas de corrección aplicadas para cada hallazgo.

\section{Remediación Específica por Vulnerabilidad Identificada}

% -----------------------------------------------------------------------------
\subsection{Remediación VULN-01: Exclusión de Suites Anónimas Diffie-Hellman en Postfix}
% -----------------------------------------------------------------------------

Para erradicar la vulnerabilidad de Man-in-the-Middle (MitM) en los puertos \texttt{25/tcp} (SMTP) y \texttt{587/tcp} (Submission):

\begin{enumerate}
  \item \textbf{Prohibición Estricta de Cifrados Anónimos (\textit{aNULL}):} En la configuración del demonio Postfix (\texttt{services/mail/config/main.cf}), se agregaron directivas que impiden de forma tajante la negociación de suites sin autenticación de servidor:
\begin{lstlisting}[style=terminal, caption={Endurecimiento Criptografico TLS en Postfix (main.cf)}, label={lst:mitigacion-postfix}]
# Excluir suites anonimas (aNULL), suites nulas (eNULL) y algoritmos obsoletos
smtpd_tls_exclude_ciphers = aNULL, eNULL, EXPORT, DES, RC4, MD5, PSK, aECDH, EDH-DSS-DES-CBC3-SHA, EDH-RSA-DES-CBC3-SHA
smtpd_tls_ciphers = high
smtpd_tls_mandatory_ciphers = high
smtpd_tls_protocols = !SSLv2, !SSLv3, !TLSv1, !TLSv1.1
smtpd_tls_mandatory_protocols = !SSLv2, !SSLv3, !TLSv1, !TLSv1.1

# Parametros Diffie-Hellman robustos de 2048/3072 bits
smtpd_tls_dh1024_param_file = /etc/postfix/dh2048.pem
\end{lstlisting}

  \item \textbf{Generación de Parámetros Diffie-Hellman Seguros:} Se sustituyeron los grupos débiles de 1024 bits generando un grupo de parámetros criptográficos de 2048 bits:
\begin{lstlisting}[style=terminal, caption={Comando OpenSSL para Parametros DH Fuertes}, label={lst:dhparam-cmd}]
openssl dhparam -out services/mail/config/dh2048.pem 2048
\end{lstlisting}
\end{enumerate}

% -----------------------------------------------------------------------------
\subsection{Remediación VULN-02: Mitigación de Denegación de Servicio Byte-Range (CVE-2011-3192)}
% -----------------------------------------------------------------------------

Para mitigar el riesgo de agotamiento de memoria por solicitudes excesivas de rangos superpuestos (\textit{overlapping byte-ranges}) en los puertos web \texttt{80/tcp} y \texttt{443/tcp}:

\begin{enumerate}
  \item \textbf{Limitación de Rangos en Nginx (\texttt{max\_ranges}):} En el archivo de configuración del proxy web (\texttt{services/web/nginx.conf}), se restringió el número máximo de rangos permitidos por solicitud:
\begin{lstlisting}[style=terminal, caption={Proteccion contra Ataques Byte-Range DoS en Nginx}, label={lst:mitigacion-ranges}]
# Limitar el procesamiento a un solo rango por peticion HTTP
max_ranges 1;

# Limitar tamano maximo del cuerpo para mitigar buffer exhausts
client_max_body_size 5m;
\end{lstlisting}

  \item \textbf{Rate Limiting en Capa 7:} Se configuraron zonas de control de concurrencia y tasa de peticiones mediante \texttt{limit\_req} y \texttt{limit\_conn}:
\begin{lstlisting}[style=terminal, caption={Limitacion de Conexiones y Tasa de Peticiones en Nginx}, label={lst:mitigacion-ratelimit}]
limit_req_zone $binary_remote_addr zone=banco_limit:10m rate=20r/s;
limit_conn_zone $binary_remote_addr zone=banco_conn:10m;

# En el bloque server
limit_conn banco_conn 20;
limit_req zone=banco_limit burst=30 nodelay;
limit_req_status 429;
\end{lstlisting}
\end{enumerate}

% -----------------------------------------------------------------------------
\subsection{Remediación VULN-03: Supresión de Métodos DEBUG y Ocultamiento de Banners}
% -----------------------------------------------------------------------------

Para resolver la advertencia de \texttt{DEBUG is enabled} en el puerto \texttt{8080/tcp} (Roundcube Webmail):

\begin{enumerate}
  \item \textbf{Desactivación de Depuración en la Aplicación:} En \texttt{services/webmail/config/config.inc.php}:
\begin{lstlisting}[style=terminal, caption={Desactivacion de Modo Debug en Roundcube Webmail}, label={lst:mitigacion-debug}]
$config['debug_level'] = 1;
$config['devel_mode'] = false;
$config['display_errors'] = false;
\end{lstlisting}

  \item \textbf{Filtrado de Métodos HTTP Anómalos en Nginx:} Se bloquea explícitamente cualquier solicitud con verbos de depuración (\texttt{DEBUG}, \texttt{TRACE}, \texttt{TRACK}):
\begin{lstlisting}[style=terminal, caption={Bloqueo de Metodos HTTP de Depuracion en Nginx}, label={lst:mitigacion-verbos}]
if ($request_method ~ ^(DEBUG|TRACE|TRACK)$) {
    return 405 "Metodo no permitido por politica de seguridad bancaria.\n";
}
server_tokens off;
\end{lstlisting}
\end{enumerate}

% -----------------------------------------------------------------------------
\subsection{Remediación VULN-04: Aislamiento del Motor de Base de Datos PostgreSQL (5432)}
% -----------------------------------------------------------------------------

Para mitigar la exposición externa del puerto de PostgreSQL:

\begin{enumerate}
  \item \textbf{Aislamiento de Red en Docker:} En el archivo \texttt{docker-compose.yml}, el puerto \texttt{5432} no se publica de forma indiscriminada a la interfaz pública, enlazándose exclusivamente a la interfaz de bucle local (\texttt{127.0.0.1:5432:5432}) o manteniéndose interno dentro de la red privada \texttt{tso-net}.
  \item \textbf{Control de Acceso Estricto en el Firewall (\texttt{nftables}):} En \texttt{server/nftables.conf}, se vinculó el acceso al puerto 5432 únicamente hacia las direcciones IP de los administradores autorizados:
\begin{lstlisting}[style=terminal, caption={Regla de Acceso a PostgreSQL en nftables.conf}, label={lst:mitigacion-nftables-db}]
# Permitir acceso al motor de BD EXCLUSIVAMENTE a IPs de administradores
tcp dport 5432 ip saddr @admin_ips accept comment "Postgres solo admins"
tcp dport 5432 drop comment "Bloqueo por defecto de Postgres a la red"
\end{lstlisting}
\end{enumerate}

% -----------------------------------------------------------------------------
\subsection{Remediación VULN-05: Cifrado y Canalización Segura de Aplicaciones Web}
% -----------------------------------------------------------------------------

Para eliminar la transmisión en texto claro en los puertos \texttt{3000}, \texttt{8000}, \texttt{8082} y \texttt{8085}:

\begin{enumerate}
  \item Se implementó la terminación SSL centralizada en Nginx con certificados autofirmados robustos generados automáticamente en el build de la imagen (\texttt{/etc/ssl/banco/server.crt}).
  \item Se aplicaron cabeceras HTTP de endurecimiento contra ataques web modernos (OWASP Top 10):
\begin{lstlisting}[style=terminal, caption={Cabeceras de Endurecimiento Web en Nginx}, label={lst:cabeceras-http}]
add_header X-Frame-Options "DENY" always;
add_header X-Content-Type-Options "nosniff" always;
add_header X-XSS-Protection "1; mode=block" always;
add_header Referrer-Policy "strict-origin-when-cross-origin" always;
\end{lstlisting}
\end{enumerate}

% -----------------------------------------------------------------------------
\subsection{Consolidación DEF-01: Mantenimiento del WAF Anti-SQLi y Suricata IPS}
% -----------------------------------------------------------------------------

Dado que los scripts de ataque \texttt{http-sql-injection} y \texttt{http-vuln-cve2017-1001000} fallaron exitosamente debido a la intervención activa de los sistemas de defensa:

\begin{enumerate}
  \item \textbf{Reglas WAF en Nginx:} Se conservan las matrices de coincidencia por expresiones regulares seguras en hexadecimal (\texttt{\textbackslash x27}, \texttt{\textbackslash x23}) que bloquean patrones de inyección SQL (\texttt{UNION SELECT}, \texttt{OR 1=1}) devolviendo \texttt{403 Forbidden}.
  \item \textbf{Suricata en Modo IPS (NFQUEUE):} Se mantienen activas las firmas de intrusión bancaria con regla de acción \texttt{drop}, registrando cada intento de intrusión en \texttt{/var/log/suricata/eve.json} para su visualización inmediata en la consola web EveBox (\texttt{http://192.168.100.85:5636/}).
\end{enumerate}

% -----------------------------------------------------------------------------
\subsection{Remediación y Endurecimiento de la Plataforma de Firma Digital (PAdES)}
% -----------------------------------------------------------------------------

Para asegurar la inviolabilidad de las operaciones criptográficas, el resguardo de las llaves privadas y la integridad de los documentos PDF firmados en la plataforma institucional:

\begin{enumerate}
  \item \textbf{Canalización Exclusiva por Proxy Inverso TLS (\texttt{firma.sudoers.lan}):}
  El contenedor \texttt{tso-firma} se mantiene confinado dentro de la red privada puente de Docker (\texttt{tso-net}), escuchando en el puerto interno \texttt{5000/tcp} sin exposición directa de sockets al exterior. Todo el tráfico de usuarios y clientes de auditoría se enruta obligatoriamente a través del servidor web Nginx con cifrado TLS 1.3:
\begin{lstlisting}[style=terminal, caption={Configuracion de Proxy Inverso Seguro para Firma Digital en Nginx}, label={lst:nginx-firma-conf}]
server {
    listen 443 ssl http2;
    server_name firma.sudoers.lan;

    ssl_certificate /etc/ssl/banco/server.crt;
    ssl_certificate_key /etc/ssl/banco/server.key;
    ssl_protocols TLSv1.2 TLSv1.3;

    location / {
        proxy_pass http://tso-firma:5000;
        proxy_set_header Host $host;
        proxy_set_header X-Real-IP $remote_addr;
        proxy_set_header X-Forwarded-For $proxy_add_x_forwarded_for;
        proxy_set_header X-Forwarded-Proto https;
    }
}
\end{lstlisting}

  \item \textbf{Custodia de Claves Privadas y Permisos Restringidos en Reposo:}
  Las llaves privadas de los certificados X.509 de los usuarios se almacenan en SQLite (\texttt{app.db}) cifradas mediante el esquema simétrico \textbf{Fernet} (\textbf{AES-128-CBC} con autenticación \textbf{HMAC-SHA256}). El archivo de la llave maestra y la base de datos se alojan en el volumen persistente \texttt{firma\_data} con permisos de archivo del sistema operativo \texttt{chmod 0600}, restringiendo su lectura al proceso del contenedor.

  \item \textbf{Protección contra Clickjacking y Fuga de Trazos Biométricos:}
  Dado que la aplicación incluye un lienzo interactivo HTML5 para la captura de rúbricas manuscritas, se inyectaron directivas perimétricas \texttt{X-Frame-Options "DENY"} y \texttt{Content-Security-Policy: frame-ancestors 'none'} para neutralizar de forma categórica ataques de \textit{UI Redressing} que intenten inducir la firma no consentida de transacciones bancarias.
\end{enumerate}
